\addcontentsline{toc}{chapter}{Executive Summary}\vspace{-1cm}
%Border
\begin{tikzpicture}[remember picture, overlay]
	\draw[line width = 4pt] ($(current page.north west) + (0.75in,-0.75in)$) rectangle ($(current page.south east) + (-0.75in,0.75in)$);
\end{tikzpicture}

The objective of this report is to provide a comprehensive account of the internship undertaken at \textbf{Hyperface} in the \textbf{Order Management System (OMS)} team. This internship served as a bridge between academic theory and industrial application, offering a profound understanding of Javascript, React, Node.js, and Redis within a professional corporate environment.

\textbf{Hyperface} is a modern fintech company that provides a credit card and payments infrastructure platform, enabling businesses to launch and manage their own branded card programs. The OMS team plays a critical role in managing the lifecycle of card orders, ensuring seamless coordination between the front-end interfaces and the backend services that power the platform.

The first phase of the internship was dedicated to understanding the organizational ethos of Hyperface, exploring its business model, product portfolio, and professional standards. This involved studying the existing codebase, understanding the microservices architecture, and familiarizing myself with the development workflow, including version control practices using Git and the agile sprint methodology followed by the team.

During the technical tenure, I was integrated into the OMS team, where I was assigned tasks related to software debugging, feature development, and performance optimization. This phase provided hands-on exposure to industry-standard tools and technologies such as \textbf{JavaScript}, \textbf{React} for building dynamic front-end interfaces, \textbf{Node.js} for server-side development, and \textbf{Redis} for efficient caching and session management. Working within a real production environment allowed me to understand the importance of writing clean, maintainable, and scalable code.

A significant milestone of the internship was the successful development and deployment of a \textbf{Progressive Web Application (PWA)}, which enhanced the accessibility and offline capabilities of the platform. Additionally, I contributed to the implementation of a \textbf{versioning system} that improved the manageability of software releases and ensured backward compatibility across different client environments. These contributions were reviewed, tested, and merged into the production codebase following standard code review processes.

